% GENERATED FILE. Edit canonical lessons, then run npm run generate:books.
% canonical-source-hash: fnv1a64:f46b23264b4a9e6e
% canonical-lessons: ZH-C432-shao, ZH-C432-danji, ZH-C432-jiejiari, ZH-C432-quanpiao, ZH-C432-guanbi

\chapter{Busy and Quiet}
\label{ch:zh-busy-and-quiet}
\hlchaptermodality{432}

\begin{chapteropening}
\textbf{By the end of this chapter:} \emph{I can say there are few visitors, talk about the off season and public holidays, buy a full-price ticket and hear that a place is closed.}

Complete the last lesson of chapter 432: I can say there are few visitors, talk about the off season and public holidays, buy a full-price ticket and hear that a place is closed.
\end{chapteropening}

\section[\texorpdfstring{\zh{少} (shǎo)}{少 (shǎo)}]{\zh{少} (shǎo) — few; not much}

\label{lesson:ZH-C432-shao}



\begin{quote}
Before the new one: say the Mandarin for many, then the Mandarin for to pose for a group photo.
\end{quote}

\subsection*{You'll want to know: \zh{少}}
\textbf{\zh{少}} — \emph{shǎo} — \textquotedblleft{}few; not much\textquotedblright{}.

\begin{quote}
\textbf{\zh{今年这家酒店的客人很少}} \emph{jīnnián zhè jiā jiǔdiàn de kèrén hěn shǎo} — This hotel has had very few guests this year.
\end{quote}

\begin{cousinweb}
\textbf{\zh{少}} (\emph{shǎo}) is inherited from Middle Chinese \emph{syewX} (Guangyun, in Baxter's transcription), glossed 'not many'. The dictionary lists a second reading, \emph{syewH}, glossed 'young': the shào of \zh{少年}.
\end{cousinweb}

\begin{grammarlens}[title={What you've built}]
One more word for getting around.
\end{grammarlens}

\subsection*{Your turn}
\begin{itemize}
\raggedright
  \item \emph{Say it:} \emph{shǎo}
  \item \emph{Say it:} \emph{shǎo}, once more
  \item \emph{Say it:} \emph{jīnnián zhè jiā jiǔdiàn de kèrén hěn shǎo} — This hotel has had very few guests this year.
  \item \emph{Recall:} say the Mandarin for many, then the Mandarin for to pose for a group photo, then say \emph{shǎo} again
\end{itemize}

\begin{tcolorbox}[breakable,colback=teal!4,colframe=teal!35!black,title={Before you move on}]
What does \zh{少} mean? (\textquotedblleft{}Few; not much\textquotedblright{}.) Say it once more.
\end{tcolorbox}
\section[dànjì]{dànjì — off season; quiet season}

\label{lesson:ZH-C432-danji}



\begin{quote}
Before the new one: say the Mandarin for to pose for a group photo, then the Mandarin for few.
\end{quote}

\subsection*{You'll want to know: dànjì}
\emph{dànjì} — \textquotedblleft{}off season; quiet season\textquotedblright{}.

\begin{quote}
\emph{dànjì de shíhou yóukè hěn shǎo} — In the off season there are very few visitors.
\end{quote}

Its characters come later; for now it is a word for the ear.

\begin{cousinweb}
\emph{dànjì} is built from \emph{dàn} 'thin, slack' and \emph{jì} 'a season', both inherited from Middle Chinese, \emph{damX} and \emph{kjwijH} (Guangyun, in Baxter's transcription), which glosses the first as a thin taste and the second as the youngest of brothers; the sense 'a season' is not given there.
\end{cousinweb}

\begin{grammarlens}[title={What you've built}]
One more word for getting around.
\end{grammarlens}

\subsection*{Your turn}
\begin{itemize}
\raggedright
  \item \emph{Say it:} \emph{dànjì}
  \item \emph{Say it:} \emph{dànjì}, once more
  \item \emph{Say it:} \emph{dànjì de shíhou yóukè hěn shǎo} — In the off season there are very few visitors.
  \item \emph{Recall:} say the Mandarin for to pose for a group photo, then the Mandarin for few, then say \emph{dànjì} again
\end{itemize}

\begin{tcolorbox}[breakable,colback=teal!4,colframe=teal!35!black,title={Before you move on}]
What does dànjì mean? (\textquotedblleft{}Off season; quiet season\textquotedblright{}.) Say it once more.
\end{tcolorbox}
\section[\texorpdfstring{\zh{节假日} (jiéjiàrì)}{节假日 (jiéjiàrì)}]{\zh{节假日} (jiéjiàrì) — public holiday}

\label{lesson:ZH-C432-jiejiari}



\begin{quote}
Before the new one: say the Mandarin for few, then the Mandarin for the off season.
\end{quote}

\subsection*{You'll want to know: \zh{节假日}}
\textbf{\zh{节假日}} — \emph{jiéjiàrì} — \textquotedblleft{}public holiday\textquotedblright{}.

\begin{quote}
\emph{jiéjiàrì de shíhou hěn duō rén zuò huǒchē qù hǎibiān} — On public holidays lots of people take the train to the seaside.
\end{quote}

\begin{cousinweb}
\textbf{\zh{节假日}} (\emph{jiéjiàrì}) is built from \zh{节} 'a festival', \zh{假} 'leave' and \zh{日} 'a day', all inherited from Middle Chinese, \emph{tset}, \emph{kaeH} and \emph{nyit} (Guangyun, in Baxter's transcription), which glosses the first as the joint of a bamboo stem and the second 'to borrow' and 'leave'; the sense 'a festival' for the first is not given there.
\end{cousinweb}

\begin{grammarlens}[title={What you've built}]
One more word for getting around.
\end{grammarlens}

\subsection*{Your turn}
\begin{itemize}
\raggedright
  \item \emph{Say it:} \emph{jiéjiàrì}
  \item \emph{Say it:} \emph{jiéjiàrì}, once more
  \item \emph{Say it:} \emph{jiéjiàrì de shíhou hěn duō rén zuò huǒchē qù hǎibiān} — On public holidays lots of people take the train to the seaside.
  \item \emph{Recall:} say the Mandarin for few, then the Mandarin for the off season, then say \emph{jiéjiàrì} again
\end{itemize}

\begin{tcolorbox}[breakable,colback=teal!4,colframe=teal!35!black,title={Before you move on}]
What does \zh{节假日} mean? (\textquotedblleft{}Public holiday\textquotedblright{}.) Say it once more.
\end{tcolorbox}
\section[\texorpdfstring{\zh{全票} (quánpiào)}{全票 (quánpiào)}]{\zh{全票} (quánpiào) — full-price ticket}

\label{lesson:ZH-C432-quanpiao}



\begin{quote}
Before the new one: say the Mandarin for the off season, then the Mandarin for a public holiday.
\end{quote}

\subsection*{You'll want to know: \zh{全票}}
\textbf{\zh{全票}} — \emph{quánpiào} — \textquotedblleft{}full-price ticket\textquotedblright{}.

\begin{quote}
\emph{wǒmen mǎi le liǎng zhāng quánpiào} — We bought two full-price tickets.
\end{quote}

\begin{cousinweb}
\textbf{\zh{全票}} (\emph{quánpiào}) is built from \zh{全} 'whole' and \zh{票} 'a ticket'. The first is inherited from Middle Chinese \emph{dzjwen} (Guangyun, in Baxter's transcription), glossed 'complete'; no dictionary we use records an older form of \zh{票}, so its history is not settled.
\end{cousinweb}

\begin{grammarlens}[title={What you've built}]
One more word for getting around.
\end{grammarlens}

\subsection*{Your turn}
\begin{itemize}
\raggedright
  \item \emph{Say it:} \emph{quánpiào}
  \item \emph{Say it:} \emph{quánpiào}, once more
  \item \emph{Say it:} \emph{wǒmen mǎi le liǎng zhāng quánpiào} — We bought two full-price tickets.
  \item \emph{Recall:} say the Mandarin for the off season, then the Mandarin for a public holiday, then say \emph{quánpiào} again
\end{itemize}

\begin{tcolorbox}[breakable,colback=teal!4,colframe=teal!35!black,title={Before you move on}]
What does \zh{全票} mean? (\textquotedblleft{}Full-price ticket\textquotedblright{}.) Say it once more.
\end{tcolorbox}
\section[guānbì]{guānbì — to shut down; to be shut to the public}

\label{lesson:ZH-C432-guanbi}



\begin{quote}
Before the new one: say the Mandarin for a public holiday, then the Mandarin for a full-price ticket.
\end{quote}

\subsection*{You'll want to know: guānbì}
\emph{guānbì} — \textquotedblleft{}to shut down; to be shut to the public\textquotedblright{}.

\begin{quote}
\emph{dōngtiān zhège jǐngqū guānbì} — This scenic area is closed in winter.
\end{quote}

Its characters come later; for now it is a word for the ear.

\begin{cousinweb}
\emph{guānbì} is built from \zh{关} 'to bar' and \emph{bì} 'to shut', both inherited from Middle Chinese, \emph{kwaen} and \emph{pejH} (Guangyun, in Baxter's transcription), which glosses the first as the bar across a door and the second 'to shut a door'.
\end{cousinweb}

\begin{grammarlens}[title={What you've built}]
One more word for getting around.
\end{grammarlens}

\subsection*{Your turn}
\begin{itemize}
\raggedright
  \item \emph{Say it:} \emph{guānbì}
  \item \emph{Say it:} \emph{guānbì}, once more
  \item \emph{Say it:} \emph{dōngtiān zhège jǐngqū guānbì} — This scenic area is closed in winter.
  \item \emph{Recall:} say the Mandarin for a public holiday, then the Mandarin for a full-price ticket, then say \emph{guānbì} again
\end{itemize}

\begin{tcolorbox}[breakable,colback=teal!4,colframe=teal!35!black,title={Before you move on}]
What does guānbì mean? (\textquotedblleft{}To shut down; to be shut to the public\textquotedblright{}.) Say it once more.
\end{tcolorbox}
